\documentclass[10pt,a4paper]{amsart}
\usepackage{amsmath,amssymb,amsthm,graphicx}
% Retained minimal preamble for the Stacks Project chapter more-algebra.tex (More on Algebra).
% The chapter file pinned-original-source.tex carries only \input{preamble} on its first line, so no
% preamble text lives inside the chapter itself.  The Stacks Project uses one single shared preamble for
% every chapter; the snapshot stages it identically as sources/stacks-chow/preamble.tex and
% sources/stacks-etale-cohomology/preamble.tex.  The block below is copied mechanically, line for line,
% from sources/stacks-chow/preamble.tex lines 14-18 (the Xy-pic setup) and 154-217 (the native theorem
% declarations together with the mathematical macros).  Nothing was added, removed or redefined; the
% wrapper class and the amsmath/amssymb/amsthm/graphicx base packages are supplied by the assembler.
% Chapter: more-algebra.tex -- More on Algebra.
% For commutative diagrams we use Xy-pic
\usepackage[all]{xy}

% We use 2cell for 2-commutative diagrams.
\xyoption{2cell}

% Theorem environments.
%
\theoremstyle{plain}
\newtheorem{theorem}[subsection]{Theorem}
\newtheorem{proposition}[subsection]{Proposition}
\newtheorem{lemma}[subsection]{Lemma}

\theoremstyle{definition}
\newtheorem{definition}[subsection]{Definition}
\newtheorem{example}[subsection]{Example}
\newtheorem{exercise}[subsection]{Exercise}
\newtheorem{situation}[subsection]{Situation}

\theoremstyle{remark}
\newtheorem{remark}[subsection]{Remark}
\newtheorem{remarks}[subsection]{Remarks}

\numberwithin{equation}{subsection}

% Macros
%
\def\lim{\mathop{\mathrm{lim}}\nolimits}
\def\colim{\mathop{\mathrm{colim}}\nolimits}
\def\Spec{\mathop{\mathrm{Spec}}}
\def\Hom{\mathop{\mathrm{Hom}}\nolimits}
\def\Ext{\mathop{\mathrm{Ext}}\nolimits}
\def\SheafHom{\mathop{\mathcal{H}\!\mathit{om}}\nolimits}
\def\SheafExt{\mathop{\mathcal{E}\!\mathit{xt}}\nolimits}
\def\Sch{\mathit{Sch}}
\def\Mor{\mathop{\mathrm{Mor}}\nolimits}
\def\Ob{\mathop{\mathrm{Ob}}\nolimits}
\def\Sh{\mathop{\mathit{Sh}}\nolimits}
\def\NL{\mathop{N\!L}\nolimits}
\def\CH{\mathop{\mathrm{CH}}\nolimits}
\def\proetale{{pro\text{-}\acute{e}tale}}
\def\etale{{\acute{e}tale}}
\def\QCoh{\mathit{QCoh}}
\def\Ker{\mathop{\mathrm{Ker}}}
\def\Im{\mathop{\mathrm{Im}}}
\def\Coker{\mathop{\mathrm{Coker}}}
\def\Coim{\mathop{\mathrm{Coim}}}

% Boxtimes
%
\DeclareMathSymbol{\boxtimes}{\mathbin}{AMSa}{"02}

%
% Macros for moduli stacks/spaces
%
\def\QCohstack{\mathcal{QC}\!\mathit{oh}}
\def\Cohstack{\mathcal{C}\!\mathit{oh}}
\def\Spacesstack{\mathcal{S}\!\mathit{paces}}
\def\Quotfunctor{\mathrm{Quot}}
\def\Hilbfunctor{\mathrm{Hilb}}
\def\Curvesstack{\mathcal{C}\!\mathit{urves}}
\def\Polarizedstack{\mathcal{P}\!\mathit{olarized}}
\def\Complexesstack{\mathcal{C}\!\mathit{omplexes}}
% \Pic is the operator that assigns to X its picard group, usage \Pic(X)
% \Picardstack_{X/B} denotes the Picard stack of X over B
% \Picardfunctor_{X/B} denotes the Picard functor of X over B
\def\Pic{\mathop{\mathrm{Pic}}\nolimits}
\def\Picardstack{\mathcal{P}\!\mathit{ic}}
\def\Picardfunctor{\mathrm{Pic}}
\def\Deformationcategory{\mathcal{D}\!\mathit{ef}}
\usepackage{cleveref}
\crefname{originallemma}{original lemma}{original lemmas}
\Crefname{originallemma}{Original lemma}{Original lemmas}
\makeatletter
\newlabel{more-algebra-lemma-check-torsion}{{0AUT}{source}{}{sourcepointer-1}{}}
\newlabel{more-algebra-lemma-check-torsion@cref}{{[originallemma][0][]0AUT}{[1][0][]source}{}{}{}}
\makeatother
\crefname{originallemma}{original lemma}{original lemmas}
\Crefname{originallemma}{Original lemma}{Original lemmas}
\makeatletter
\newlabel{lemma-localize-torsion}{{0AUR}{source}{}{sourcepointer-2}{}}
\newlabel{lemma-localize-torsion@cref}{{[originallemma][0][]0AUR}{[1][0][]source}{}{}{}}
\makeatother
\crefname{originallemma}{original lemma}{original lemmas}
\Crefname{originallemma}{Original lemma}{Original lemmas}
\makeatletter
\newlabel{algebra-lemma-characterize-zero-local}{{00HN}{source}{}{sourcepointer-3}{}}
\newlabel{algebra-lemma-characterize-zero-local@cref}{{[originallemma][0][]00HN}{[1][0][]source}{}{}{}}
\makeatother
\title{Torsion freeness is detected on localisations at maximal ideals}
\author{The Stacks Project Authors}
\date{stacks-project:more-algebra.tex:local-snapshot-20261012; GNU FDL 1.2 or later}
\begin{document}
\maketitle
\setcounter{subsection}{1}
\noindent\textit{Editorial note.} All delivered mathematical text is the author's own, copied from the pinned Stacks Project chapter file more-algebra.tex; no statement or proof step was rewritten, reordered or removed.\par\medskip
\noindent\textit{Editorial note.} The chapter file pinned-original-source.tex carries only \textbackslash{}input\{preamble\} on its first line, so no preamble text lives inside the chapter. The retained file sources/stacks-more-algebra/retained-preamble.tex is a mechanical, line-for-line copy of the project's single shared preamble as staged in the snapshot, namely sources/stacks-chow/preamble.tex lines 14-18 (Xy-pic setup) and 154-217 (theorem declarations and mathematical macros). Nothing was added, removed or redefined, no package was added, and the excerpt serves only as preamble, never as a span coordinate.\par\medskip
\noindent\textit{Editorial note.} Both counted spans start at the \textbackslash{}begin line of their environment, so the statement and the author's proof are delivered whole; surrounding prose of the section is deliberately outside the counted spans and is carried as standard prerequisites instead.\par\medskip
\noindent\textit{Editorial note.} The citation fields (source\_id, source\_version, authors, source\_url, license.name, source\_title) are identical to those of the cards already published from this chapter, so that the batch selector treats them as one source instead of deferring them as a same-source clash.\par\medskip
\noindent\textit{Allowed standard prerequisites (editorial).}\begin{itemize}
\item Torsion element and torsion free module over a domain. definition at the start of the section, source lines 4781-4792. Supplies the condition that is being tested locally.
\item A localisation of a torsion free module is torsion free. the statement counted by the sibling card S1-algebra-localize-torsion at 4813-4822, cited in the hint at proof lines 4859-4860. Gives the direction from torsion freeness of M to torsion freeness of its localisations.
\item A module is zero if and only if all its localisations at maximal ideals are zero. the author cites Algebra, Lemma algebra-lemma-characterize-zero-local in the hint at proof line 4860. Converts the local vanishing of the annihilator of a putative torsion element into its vanishing over the ring.
\end{itemize}\medskip
\ifdefined\hypertarget\hypertarget{sourcepointer-1}{}\fi
\noindent\textit{Source reference.} Original lemma 0AUT: focused-seeds/AGP/sources/stacks-snapshot-20261011/tags.txt, source lines 13661--13661. The lemma's own canonical tag, quoted so that readers can resolve the label inside the delivered text; it points at the counted statement itself and carries no proof content.\par\medskip
\ifdefined\hypertarget\hypertarget{sourcepointer-2}{}\fi
\noindent\textit{Source reference.} Original lemma 0AUR: focused-seeds/AGP/sources/stacks-snapshot-20261011/tags.txt, source lines 13659--13659. Navigation pointer to the statement of the same section that a localisation of a torsion free module is torsion free, which the author's hint invokes; nothing here re-proves it.\par\medskip
\ifdefined\hypertarget\hypertarget{sourcepointer-3}{}\fi
\noindent\textit{Source reference.} Original lemma 00HN: focused-seeds/AGP/sources/stacks-snapshot-20261011/tags.txt, source lines 649--649. Navigation pointer into the Algebra chapter for the criterion that a module vanishes when all its localisations at maximal ideals do; the author names it in the hint and it is a standard input.\par\medskip
\begin{lemma}
\label{lemma-check-torsion}
Let $R$ be a domain. Let $M$ be an $R$-module.
Then $M$ is torsion free if and only if $M_\mathfrak m$ is a
torsion free $R_\mathfrak m$-module for all maximal ideals
$\mathfrak m$ of $R$.
\end{lemma}

\begin{proof}
Omitted. Hint: Use Lemma \ref{lemma-localize-torsion} and
Algebra, Lemma \ref{algebra-lemma-characterize-zero-local}.
\end{proof}

\end{document}
